% =============================================================================
\section{Additive-Model Derivations and Variance Identities}\label{app:additive}
% =============================================================================

\subsection{Proof of \texorpdfstring{\cref{thm:sobol_equiv}}{Theorem: Sobol Equivalence}}

\begin{proof}
Under the additive model $f(X) = \sum_{i=1}^L g_i(X_i)$ with independent coordinates and independent-copy perturbation at position $i$:
\[
f(X) - f(X^{(i)}) = g_i(X_i) - g_i(X_i'),
\]
where $X_i' \sim P_i$ is independent of $X_i$. Therefore:
\[
\cI(i;r) = \E\!\left[\norm{g_i(X_i) - g_i(X_i')}_2^2\right] = 2\Var(g_i(X_i)),
\]
using the identity $\E[(Y-Y')^2] = 2\Var(Y)$ for i.i.d.\ copies $Y, Y'$.

For a scalar projection $\varphi(f(X)) = \sum_i \varphi(g_i(X_i))$, the total variance decomposes as $V_{\mathrm{tot}} = \sum_i \Var(\varphi(g_i(X_i)))$ by independence. The Sobol total-effect index for coordinate $i$ in an additive model with independent inputs equals the first-order index: $S_i^T = S_i = \Var(\varphi(g_i(X_i)))/V_{\mathrm{tot}}$. Therefore $\cI(i;r) = 2V_{\mathrm{tot}} S_i^T$.
\end{proof}

\subsection{Quantitative tail approximation (\texorpdfstring{\cref{prop:tail_add}}{Proposition: Tail Approximation})}

\begin{proof}
Let $f(X)=\sum_{i=1}^L g_i(X_i) + \epsilon(X)$ and $S = \{i:|i-r|>\ell\}$. Define:
\begin{align*}
A_S &:= \sum_{i\in S}(g_i(X_i)-g_i(X_i')), \\
B_S &:= \epsilon(X)-\epsilon(X^{(S)}).
\end{align*}
Then $\Delta_S := f(X) - f(X^{(S)}) = A_S + B_S$ and:
\[
\E\norm{\Delta_S}^2 = \E\norm{A_S}^2 + 2\E\ip{A_S}{B_S} + \E\norm{B_S}^2.
\]
Under independence and the additive structure:
\[
\E\norm{A_S}^2 = \sum_{i\in S}\E\norm{g_i(X_i)-g_i(X_i')}^2 = \sum_{i\in S} \cI(i;r).
\]
By Cauchy--Schwarz: $|2\E\ip{A_S}{B_S}| \le 2\sqrt{\E\norm{A_S}^2}\sqrt{\E\norm{B_S}^2}$.

Since $\E\norm{B_S}^2 \le 4\eta^2$ (where $\eta^2 = \E[\norm{\epsilon(X)}^2]$, using triangle inequality and i.i.d.\ structure):
\[
\left|\Delta_{\mathrm{out}}(\ell;r) - \sum_{i\in S}\cI(i;r)\right| \le 2\sqrt{\sum_{i\in S}\cI(i;r)} \cdot 2\eta + 4\eta^2.
\]
\end{proof}

\subsection{Variance decomposition identity}

For completeness, we record the law of total variance:
\[
\Var(Y) = \E[\Var(Y\mid \cG)] + \Var(\E[Y\mid \cG]),
\]
where $\cG$ is a sub-$\sigma$-algebra. This underlies the Sobol decomposition: conditioning on $X_S$ yields $V_S := \Var(\E[Y\mid X_S])$, the variance explained by coordinates in $S$.

\subsection{Explicit computation for quadratic models}

Consider $f(X) = \sum_i g_i(X_i) + \sum_{i<j} g_{ij}(X_i, X_j)$ (pairwise interactions). The influence energy for perturbing position $i$ becomes:
\[
\cI(i;r) = 2\Var(g_i(X_i)) + 2\sum_{j \neq i} \E\!\left[\Var(g_{ij}(X_i, X_j) \mid X_j)\right],
\]
where the second term captures the interaction contribution. This shows that single-site perturbation captures both main effects and the ``average interaction'' with all other positions, but not the specific pairwise interaction structure.

\subsection{Whitening invariance}

Let $\Sigma=\Cov(f(X))$ be positive definite and define whitened embeddings $\tilde f(X)=\Sigma^{-1/2}(f(X)-\E[f(X)])$. For any invertible linear transform $T$, the model $Tf$ has covariance $T\Sigma T^\top$. Whitening $(T\Sigma T^\top)^{-1/2}(Tf(X) - T\E[f(X)]) = (T\Sigma T^\top)^{-1/2}T\Sigma^{1/2}\tilde{f}(X)$. Since $(T\Sigma T^\top)^{-1/2}T\Sigma^{1/2}$ is orthogonal (verifiable by checking that its square is the identity up to sign), ECL under Mahalanobis distance is invariant under invertible linear reparameterization.
